\documentclass[10pt]{article}
\usepackage[margin=0.82in]{geometry}
\usepackage{amsmath,amssymb,mathtools}
\usepackage[T1]{fontenc}
\usepackage[hidelinks]{hyperref}
\usepackage{microtype}

\newcommand{\D}{\mathcal D}
\newcommand{\W}{\mathcal W}
\newcommand{\B}{\mathbb B}

\title{The Higher-Dimensional Spin-Ball/Max-Ball\
Question Has a Negative Answer}
\author{fa\_banach\_001 / agent\_lane\_00}
\date{17 August 2026}

\begin{document}
\maketitle

\noindent\textbf{Status: literature already answered (full answer for every
$m\geq 3$).}

\section*{Original question}

Farenick, Huntinghawk, Masanika, and Plosker prove that their spin ball and
max ball agree in dimensions one and two.  Immediately after
Theorem~5.12, on PDF page~19, they ask whether the theorem extends to higher
dimensions \cite{FHMP}.

The corrected universal three-variable coefficient tuple is
\[
 F^{[3]}\cong P\oplus\overline P,
 \qquad
 \D_{F^{[3]}}=\D_P\cap\D_{\overline P}.
\]
This point matters: testing only the irreducible Pauli triple $P$ produces a
spurious level-two ``counterexample'' that merely detects the failure of the
transpose map to be completely positive.  It is excluded by the conjugate
summand in the corrected universal tuple.

\section*{Explicit later answer}

Evert and Passer restate the same problem as Question~1 on PDF page~6,
explicitly citing the text after Farenick et al.'s Theorem~5.12
\cite{EP}.  Their Theorem~2.9 (PDF page~10) constructs free extreme points of
$\D_{F^{[3]}}$ at matrix levels four and six.  Hence
$\D_{F^{[3]}}$ is not the minimal matrix convex set over the Euclidean ball,
and its free polar dual is not maximal.  Remark~2.11 (PDF page~11) states
directly that this gives a negative answer to the spin-ball/max-ball
equality question in three or more variables (accounting for the reversed
min/max terminology in the source paper).

Finally, Corollary~3.9 (PDF page~19) proves for every $g\geq3$ that
$\D_{F^{[g]}}$ has a higher-level free extreme point and therefore is not the
minimal matrix convex set over $\B_g$.  Thus the proposed equality fails in
every higher dimension, not merely for $g=3$.

The supporting authors knew they were resolving the source question.  This
is therefore an explicit literature answer rather than an agent-inferred
reformulation.

\section*{Scope and search evidence}

The source question is fully resolved: equality holds for $m=1,2$ and fails
for every $m\geq3$.  The later paper was found by checking the citation graph
of arXiv:2006.06810 and then matching its Question~1 and Remark~2.11 against
the exact source location.  No new proof or counterexample is claimed here.

\begin{thebibliography}{9}

\bibitem{FHMP}
D. Farenick, F. Huntinghawk, A. Masanika, and S. Plosker,
\emph{Complete order equivalence of spin unitaries},
Linear Algebra Appl. \textbf{610} (2021), 1--28;
arXiv:2006.06810.  Corrigendum: Linear Algebra Appl. \textbf{635}
(2022), 201--202.

\bibitem{EP}
E. Evert and B. Passer,
\emph{Matrix convex sets over the Euclidean ball and polar duals of real
free spectrahedra}, arXiv:2507.20325 (2025),
doi:10.1016/j.laa.2026.01.015.

\end{thebibliography}

\end{document}
